\chapter{Cardinal Characteristics of the Continuum}

\section{Growth of Functions}

\label{b,d-ineq}
\textcolor{laser_eyes}{Definição de $\mf{b}$ e $\mf{d}$ e demonstração de que
$$\aleph_1\leq\text{cf}(\mf{b})=\mf{b}\leq\text{cf}(\mf{d})\leq\mf{d}\leq\mf{c}$$}
\textcolor{laser_eyes}{(Andreas Blass - Handbook of Set Theory [6])}

\section{Galois-Tukey Connections}

Lets introduce the language of Galois Tukey Connections (GT).

\begin{definition}{}{}
    A relational system is a triple $\mc{A}=(A_-,A_+,A)$, with $A\subseteq A_-\times A_+$ where, for every $a_-\in A_-$, there is $a_+\in A_+$ such that $a_-Aa_+$, and, for every $a_+\in A_+$ there is $a_-\in A_-$ such that $a_-Aa_+$ fails.
\end{definition}

You can interpret $A_-$ as the set of challenges you need to surpass, and $A_+$ as the responses you will use to pass the challenges. So, the relation $A$ indicates when a challenge is passed. The other conditions requires that any challenge can be surpassed (this guarantees $\mf{d}(\mc{A})$ is well defined), and that we don't have any response for every challenge, to avoid trivializations (and to guarantee $\mf{b}(\mc{A})$ is well defined).

Any relational system $\mc{A}=(A_-,A_+,A)$ has a dual $A^\top=(A_+,A_-,\neg\check{A})$, where $a_+\neg\check{A}a_-$ if $\neg(a_-Aa_+)$.

\begin{definition}{}{}
    Given a relation system $\mc{A}=(A_-,A_+,A)$, we associate the following cardinals
    \begin{itemize}
        \item the dominating number $\mf{d}(\mc{A})$ is the smallest size of a family $C_+\subseteq A_+$ s.t. for all $a_-\in A_-$, there is a $a_+\in C_+$ with $a_-Aa_+$;
        \item the unbounding number $\mf{b}(\mc{A})$ is the smallest size of a family $C_-\subseteq A_-$ s.t. for all $a_+\in A_+$, there is $a_-\in C_-$ s.t. $\neg(a_-Aa_+)$.
    \end{itemize}
\end{definition}

Intuitively, the dominating number is the smallest size of a set of responses that do the job, and the unbounding number the smallest size of a set of challenges that "breaks" any given response. Note that $\mf{b}(\mc{A})$ and $\mf{d}(\mc{A})$ are dual in the sense that $\mf{b}\rp{A^\top}=\mf{d}(A)$ and $\mf{d}\rp{A^\top}=\mf{b}(A)$.

\begin{sketch}
Now, we want to define a reduction $\varphi:\mc{A}\to\mc{B}$ between two relational systems, which helps to prove that $\mf{d}(\mc{A})\leq\mf{d}(\mc{B})$ and $\mf{b}(\mc{A})\geq\mf{d}(\mc{B})$. So, it's natural to think like this: Since we want to reduce $\mc{A}$ to $\mc{B}$, we can work really well in $\mc{B}$, so, given a challenge $a_-\in A_-$, we want to convert it to a challenge $\varphi_-(a_-)\in B_-$, where we know it has an answer $b_+\in B_+$ that responds this challenge, i.e., $\varphi(a_-)B b_+$, so it only remains to convert this response $b_+\in B_+$ back to a response $\varphi_+(b_+)\in A_+$, where $a_-A\varphi_+(b_+)$.
\end{sketch}

So we can define:

\begin{definition}{}{}
    Given two relational systems $\mc{A}=(A_-,A_+,A)$ and $\mc{B}=(B_-,B_+,B)$, $\mc{A}$ is Tukey reducible to $\mc{B}$ ($\mc{A}\leq_T\mc{B}$) if there are functions $\varphi_-:A_-\to B_-$, $\varphi_+:B_+\to A_+$ such that $\varphi_-(a_-)Bb_+\Rightarrow a_-A\varphi(b_+)$ for all $a_-\in A_-$ and $b_+\in B_+$.
\end{definition}

\textcolor{laser_eyes}{Pendente}

\section{Null and Meager Sets}

The main result of this chapter is known as the Cichoń's Diagram (Figure \ref{fig:cichons-diagram})

\begin{figure}[ht]
\centering
\begin{tikzcd}
	& {\text{cov}(\mathcal{N})} & {\text{non}(\mathcal{M})} & {\cof{\mathcal{M}}} & {\cof{\mathcal{N}}} & {2^{\aleph_0}} \\
	&& {\mathfrak{b}} & {\mathfrak{d}} \\
	{\aleph_0} & {\add{\mathcal{N}}} & {\add{\mathcal{M}}} & {\text{cov}(\mathcal{M})} & {\text{non}(\mathcal{N})}
	\arrow["{\hyperref[th:rothberger]{(3)}}", from=1-2, to=1-3]
	\arrow["{\hyperref[cor:M_N_ideals]{(2)}}", from=1-3, to=1-4]
	\arrow["{(?)}", from=1-4, to=1-5]
	\arrow["{\hyperref[cor:M_N_ideals]{(2)}}", from=1-5, to=1-6]
	\arrow["{(?)}", from=2-3, to=1-3]
	\arrow["{\hyperref[th:b_d_ineq]{(1)}}", from=2-3, to=2-4]
	\arrow["{(?)}"', from=2-4, to=1-4]
	\arrow["{\hyperref[cor:M_N_ideals]{(2)}}"', from=3-1, to=3-2]
	\arrow["{\hyperref[cor:M_N_ideals]{(2)}}", from=3-2, to=1-2]
	\arrow["{(?)}"', from=3-2, to=3-3]
	\arrow["{(?)}", from=3-3, to=2-3]
	\arrow["{\hyperref[cor:M_N_ideals]{(2)}}"', from=3-3, to=3-4]
	\arrow["{(?)}"', from=3-4, to=2-4]
	\arrow["{\hyperref[th:rothberger]{(3)}}"', from=3-4, to=3-5]
	\arrow["{\hyperref[cor:M_N_ideals]{(2)}}"', from=3-5, to=1-5]
\end{tikzcd}
\caption{Cichoń's Diagram (w/ clickable hyperlinks)}
\label{fig:cichons-diagram}
\end{figure}

\textcolor{laser_eyes}{Definição de um ideal}

\begin{theorem}{}{}\label{ideals-ineq}
    If $\mc{I}$ is a $\sigma$-complete ideal containing all singletons, then
    $$\aleph_0\leq\text{add}(\mc{I})\leq\text{cov}(\mc{I}),\text{non}(\mc{I})\leq\text{cof}(\mc{I})\leq2^{\aleph_0}$$
    summarized by (Figure \ref{fig:ideals-diagram})
\end{theorem}

\begin{figure}[ht]
\centering
\begin{tikzcd}
	&& {\text{cov}(\mathcal{I})} \\
	{\aleph_0} & {\text{add}(\mathcal{I})} && {\text{cof}(\mathcal{I})} & {2^{\aleph_0}} \\
	&& {\text{non}(\mathcal{I})}
	\arrow[from=1-3, to=2-4]
	\arrow[from=2-1, to=2-2]
	\arrow[from=2-2, to=1-3]
	\arrow[from=2-2, to=3-3]
	\arrow[from=2-4, to=2-5]
	\arrow[from=3-3, to=2-4]
\end{tikzcd}
\caption{Theorem \ref{ideals-ineq}}
\label{fig:ideals-diagram}
\end{figure}

\textcolor{laser_eyes}{Definição de $\mc{M}$ e $\mc{N}$}

\begin{theorem}{}{}\label{M,N-ideals}
    $\mc{M}$ and $\mc{N}$ are both ideals, in particular, Theorem \ref{ideals-ineq} applies to them
\end{theorem}

\begin{proof}
    \textcolor{laser_eyes}{Pendente}
\end{proof}

The following classical lemma shows that null and meager describe "smallness" in two different ways.

\begin{lemma}{}{}
    There exists sets $A\in\mc{N}$ and $B\in\mc{M}$ such that $A\cup B=\mbb{R}$.
\end{lemma}

\begin{proof}
    \textcolor{laser_eyes}{Pendente}
\end{proof}

This easy lemma has two interesting consequences.

\begin{theorem}{}{}(Rothberger - 1938)\label{cov<=non}
$$\text{cov}(\mc{M})\leq\text{non}(\mc{N})\quad\text{and}\quad\text{cov}(\mc{N})\leq\text{non}(\mc{M})$$  
\end{theorem}

\begin{proof}
    \textcolor{laser_eyes}{Pendente}
\end{proof}

\begin{theorem}{}{}(Erdös-Sierpiński). If $\text{add}(\mc{N})=\text{cof}(\mc{N})$, then, there exists a bijection $f:\mbb{R}\to\mbb{R}$ such that
$$f(X)\in\mc{N}\iff X\in\mc{M}\quad\text{and}\quad f(X)\in\mc{M}\iff X\in\mc{N}$$
    In particular, $\text{CH}$ implies the existence of such bijection.
\end{theorem}

\begin{proof}
    \textcolor{laser_eyes}{Pendente}
\end{proof}